\section{Two Sessions, Each Other's Speakers}
\label{sec:ch11-two-sessions-each-others-speakers}

\index{representation order!the failure it produces|(}%
A client asks the graph for two talks by their identifiers, which is the shape
of every screen that shows a reader the things they bookmarked:

\begin{minted}[fontsize=\footnotesize,breakanywhere]{text}
POST /graphql HTTP/1.1
Host: localhost:3002
Content-Type: application/json

{"query":"{ nodes(ids: [\"U2Vzc2lvbjo0\",\"U2Vzc2lvbjox\"]) { ... on Session { title speaker { name } } } }"}
\end{minted}

\begin{minted}[fontsize=\footnotesize]{json}
{"data":{"nodes":[
  {"title":"The Bank That Split Its Graph",
   "speaker":{"name":"Ada Fischer"}},
  {"title":"Schemas That Outlive Their Authors",
   "speaker":{"name":"Chidi Okafor"}}]}}
\end{minted}

\index{representation order!both names are real}%
Both of those are wrong. The bank talk is Chidi Okafor's and the schema talk is
Ada Fischer's. The graph has told a client otherwise at HTTP 200 with no error
key, and the two names it used belong to real speakers who really are in the
database. Nothing is missing and nothing is malformed, so there is no null to
notice and nothing to catch the eye, and a person reading it cannot tell.

\index{representation order!every other page still right}%
This is not a graph that is down. Every request this book has printed since
chapter~\ref{ch:data-without-n-plus-1} still answers correctly against these
same three services. The page chapter~\ref{ch:router-n-plus-1} opened on:

\begin{minted}{graphql}
{ sessions(first: 10) { nodes { title speaker { name } } } }
\end{minted}

\index{statement count!unchanged by the defect}%
comes back with Ada Fischer twice, then Bruno Kaminski, then Chidi Okafor,
against the right four titles, at one statement in the Sessions service and one
in the Speakers service. That is the number
chapter~\ref{ch:second-third-subgraph} measured and
chapter~\ref{ch:router-n-plus-1} restored, unchanged. The default page of two
is right. The graph composes. Each service exports the schema it exported
before, byte for byte. The query plan for that page is the same
\code{Sequence} of two steps with the same \code{BatchEntity} at the same field
path.

\subsection{The File That Did It}

\index{DataLoader!written out by hand}%
Two files differ from the graph chapter~\ref{ch:router-n-plus-1} left, and the
first is \code{SpeakerDataLoader.cs}. The Speakers service's DataLoader is no
longer generated from a \code{[DataLoader]} method; it is written out by hand:

\begin{minted}[highlightlines={37-42}]{csharp}
using GreenDonut;
using Microsoft.EntityFrameworkCore;

namespace Speakers;

/// <summary>
/// What the source-generated loader is underneath: a class deriving
/// from DataLoaderBase, given the batch of keys and a span to put the
/// answers in. Written out by hand because a store that is not a
/// DbContext - a remote service, a cache, a stored procedure - hands
/// back a list rather than a dictionary, and the generator only
/// accepts a dictionary.
/// </summary>
public interface ISpeakerByIdDataLoader : IDataLoader<int, Speaker>;

public sealed class SpeakerByIdDataLoader(
    IServiceProvider services,
    IBatchScheduler batchScheduler,
    DataLoaderOptions options)
    : DataLoaderBase<int, Speaker>(batchScheduler, options),
        ISpeakerByIdDataLoader
{
    protected override async ValueTask FetchAsync(
        IReadOnlyList<int> ids,
        Memory<Result<Speaker?>> results,
        DataLoaderFetchContext<Speaker> context,
        CancellationToken ct)
    {
        await using var scope = services.CreateAsyncScope();
        var db = scope.ServiceProvider
            .GetRequiredService<SpeakerContext>();

        var rows = await db.Speakers
            .Where(s => ids.Contains(s.Id))
            .ToListAsync(ct);

        // One row per key, in the order the keys were given, so the
        // i-th row is the answer to the i-th key.
        for (var i = 0; i < rows.Count; i++)
        {
            results.Span[i] = Result<Speaker?>.Resolve(rows[i]);
        }
    }
}
\end{minted}

\index{Program.cs@\code{Program.cs}!registering a loader by hand}%
The second file is \code{Program.cs}. The source generator does not know the
loader exists, so it is registered the way anything else is. Four lines go in
above the builder chain and nothing else in the file that
chapter~\ref{ch:second-third-subgraph} printed moves:

\begin{minted}[highlightlines={26-30}]{csharp}
using ChilliCream.Nitro.App;
using HotChocolate.AspNetCore;
using Microsoft.EntityFrameworkCore;
using Speakers;

var builder = WebApplication.CreateBuilder(args);

// EF Core reports every command through the logging pipeline as well,
// at Information and with a duration attached. StatementLog below is
// this service's instrument, so silence the other one rather than have
// every statement printed twice.
builder.Logging.AddFilter(
    "Microsoft.EntityFrameworkCore.Database.Command",
    LogLevel.Warning);

builder.Services.AddDbContextFactory<SpeakerContext>(options =>
{
    options.UseSqlite("Data Source=speakers.db");

    if (builder.Environment.IsDevelopment())
    {
        options.AddInterceptors(new StatementLog());
    }
});

// A hand-written loader is not found by the source generator, so it
// is registered here like any other service.
builder.Services
    .AddDataLoader<ISpeakerByIdDataLoader, SpeakerByIdDataLoader>();

builder
    .AddGraphQL()
    .AddSpeakersTypes()
    .RegisterDbContextFactory<SpeakerContext>()
    .AddQueryContext()
    .AddPagingArguments()
    .AddGlobalObjectIdentification()
    .AddApolloFederation()
    .ModifyCostOptions(options => options.ApplyCostDefaults = false)
    .TryAddTypeInterceptor<ShareNodeFields>();

var app = builder.Build();

using (var scope = app.Services.CreateScope())
{
    var factory = scope.ServiceProvider
        .GetRequiredService<IDbContextFactory<SpeakerContext>>();
    using var db = factory.CreateDbContext();
    db.Database.EnsureCreated();
}

app.MapGraphQL()
    .WithOptions(options =>
    {
        // Serve the Nitro build that shipped in the package. The
        // default, ServeMode.Latest, downloads the current one
        // from ChilliCream's CDN instead.
        options.Tool.ServeMode = ServeMode.Embedded;
    });

app.RunWithGraphQLCommands(args);
\end{minted}

\index{DataLoader!the comment that is the bug}%
The highlighted comment and the loop under it are the whole of the defect, and
the comment is the mistake stated out loud. It says the rows come back in the
order the keys were given. They do not.
\code{ids.Contains(s.Id)} becomes an \code{IN} list, and an \code{IN} list is a
membership test rather than an instruction about order. What comes back is in
whatever order the store chose, and the store was not told what the caller
wanted. Every run of this against SQLite handed the rows back by ascending
primary key, so asking for speakers three and one gets them back as one and
three.

\index{reference resolver!untouched by the defect}%
Nothing else in the service moved. \code{SpeakerType.cs} is the file
chapter~\ref{ch:second-third-subgraph} printed, taking the same
\code{ISpeakerByIdDataLoader} it always took, decoding the same key and
checking the same type name inside it before spending it. The reference
resolver asks for the speaker it was asked about. What comes back from
\code{LoadAsync} is somebody else.
\index{representation order!the failure it produces|)}%
